% standard beamer lecture template for slides
% by Derek Huang
\documentclass{beamer}
\usepackage[utf8]{inputenc}
\usepackage{algorithm2e, amsmath, amssymb, amsfonts, graphicx}
% allow section.equation numbering
\numberwithin{equation}{section}
% use boadilla theme
\usetheme{Boadilla}
% remove navigation symbols
\usenavigationsymbolstemplate{}
% get numbered figure captions
\setbeamertemplate{caption}[numbered]
% changes itemize to circle + other things
\useoutertheme{split}
\useinnertheme{circles}

% command for the title string. change for each lecture
\newcommand{\lecturetitle}{Probabilistic generative models}
% allow automatic alert-highlighted references and hyperlinks
\newcommand{\aref}[1]{\alert{\ref{#1}}}
\newcommand{\ahref}[2]{\href{#1}{\alert{#2}}}
% title page stuff. brackets content displayed in footer bar
\title[\lecturetitle]{\lecturetitle}
% metadata. content in brackets is displayed in footer bar
\author[Derek Huang (BAC Advanced Team)]{Derek Huang}
\institute{BAC Advanced Team}
\date{March 29, 2021}

% change "ball" bullet to numbered bullet and section title for section
\setbeamertemplate{section in toc}{\inserttocsectionnumber.~\inserttocsection}
% change ball to gray square (copied from stackoverflow; \par needed for break)
\setbeamertemplate{subsection in toc}{        
    \hspace{1.2em}{\color{gray}\rule[0.3ex]{3pt}{3pt}}~\inserttocsubsection\par}
% use default enumeration scheme
\setbeamertemplate{enumerate items}[default]
% required line that fixes the problem of \mathbf, \bf not working in beamer
% for later (post-2019) TeX Live installations. see the issue on GitHub:
% https://github.com/josephwright/beamer/issues/630
\DeclareFontShape{OT1}{cmss}{b}{n}{<->ssub * cmss/bx/n}{}

\begin{document}

% title slide
\begin{frame}
    \titlepage
    \centering
    % relative path may need to be updated depending on .tex file location
    \includegraphics[scale = 0.1]{../bac_logo1.png}
\end{frame}

% table of contents slide
\begin{frame}{Overview}
    \tableofcontents
\end{frame}

% section
\section{Generative modeling}

\begin{frame}{Motivation}
    \begin{itemize}
        \item
        Let $ X : \Omega \rightarrow \mathbb{R}^d $ be the input variable,
        $ Y : \Omega \rightarrow \mathbb{R} $ be the response variable. Let
        $ \mathcal{D} \triangleq \{(\mathbf{x}_1, y_1), \ldots
        (\mathbf{x}_N, y_N)\} $ be the training data, each $ (\mathbf{x}_k,
        y_k) $ sampled independently from $ X, Y $.

        \item
        Assume a parametric model for $ X, Y $ joint or $ Y \mid X $
        conditional distribution, i.e. using a distribution parametrized by
        $ \theta \in \Theta $, $ \Theta \subseteq \mathbb{R}^p $.

        \item
        We showed that the joint likelihood $ p(\mathcal{D} \mid \theta) $
        is such that
        \begin{equation*}
            p(\mathcal{D} \mid \theta) =
            \prod_{k = 1}^Np(y_k, \mathbf{x}_k \mid \theta) \propto
            \prod_{k = 1}^Np(y_k \mid \mathbf{x}_n, \theta)
        \end{equation*}

        \item
        In linear regression, we directly modeled $ p(y \mid \mathbf{x} \mid
        \theta) $, positing $ Y \mid X \sim \mathcal{N}(\mathbf{w}^\top X + b,
        \sigma^2) $, for $ \mathbf{w} \in \mathbb{R}^d $,
        $ b \in \mathbb{R} $, $ \sigma \in (0, \infty) $.

        \item
        What's a way to model $ p(y, \mathbf{x} \mid \theta) $, the joint
        $ X, Y $ likelihood?
    \end{itemize}
\end{frame}

\subsection{Conditioning on class}

\begin{frame}{Conditioning on class}
    \begin{itemize}
        \item
        We can apply Bayes' rule to $ p(\mathcal{D} \mid \theta) $ to see that
        \begin{equation*}
            p(\mathcal{D} \mid \theta) =
            \prod_{k = 1}^Np(y_k, \mathbf{x}_k \mid \theta) =
            \prod_{k = 1}^Np(\mathbf{x}_k \mid y_k, \theta)p(y_k \mid \theta)
        \end{equation*}

        \item
        Not too helpful for regression, where usually
        $ Y(\Omega) = \mathbb{R} $, but useful for classification, where
        $ Y(\Omega) = \mathcal{C} $, $ \mathcal{C} $ some finite set of
        classes.

        \item
        $ p(\mathbf{x} \mid y, \theta) $ is a \textit{class-conditional
        likelihood}, giving the likelihood of $ \mathbf{x} $ given $ \theta $
        and knowledge that the pair $ (\mathbf{x}, y) $ is in some class.

        \item
        $ p(y \mid \theta) $ is the \textit{class prior}, essentially
        $ \mathbb{P}\{Y = C \mid \theta\} $, $ C \in \mathcal{C} $.

        \item
        With $ p(\mathbf{x} \mid y, \theta) $ and $ p(y \mid \theta) $
        we can model $ p(\mathcal{D} \mid \theta) $ directly and even generate
        new $ (\mathbf{x}, y) $ examples, hence the moniker
        \textit{generative modeling}.
    \end{itemize}
\end{frame}

\begin{frame}{Conditioning on class}
    \begin{itemize}
        \item
        Generative modeling typically limited to classification, so from now on
        assume $ Y(\Omega) = \mathcal{C} $, $ \mathcal{C} \triangleq \{C_1,
        \ldots C_K\} $, $ K $ number of classes.

        \item
        We usually model $ Y \mid \theta $ with the categorical distribution,
        i.e.
        \begin{equation} \label{cat_dist_like}
            p(y \mid \theta) \triangleq
            \sum_{k = 1}^Kp_k\mathbb{I}_{\{C_k\}}(y)
        \end{equation}
        Here $ K \in \mathbb{N} $, $ p_1, \ldots p_K \in (0, 1) $,
        $ \sum_{k = 1}^Kp_k = 1 $.

        \item        
        The \textit{class priors} $ p_1, \ldots p_K $ are typically estimated
        through maximum likelihood, i.e. $ \hat{p}_j \triangleq \frac{1}{N}
        \sum_{k = 1}^N\mathbb{I}_{\{C_j\}}(y_k) $ is the estimate for $ p_j $.
    \end{itemize}
\end{frame}

\subsection{Discrete features}

\begin{frame}{Discrete features}
    \begin{itemize}
        \item
        Suppose $ X(\Omega) = \{0, 1\}^d $, where for $ i \in
        \{1, \ldots d\} $, $ j \in \{1, \ldots K\} $,
        $ p_{11}, \ldots p_{Kd} \in (0, 1) $,
        $ X_i \mid Y = C_j \sim \operatorname{Bernoulli}(p_{ji}) $,
         $ X_1 \mid Y, \ldots X_d \mid Y $ mutually independent\footnote{
            Class-conditional input feature independence is the
            \textit{naive Bayes} assumption \cite{bishop_ml}.
        }. Let $ \mathbf{p}_j \triangleq 
        [ \ p_{j1} \ \ldots  \ p_{jd} \ ]^\top \in \mathbb{R}^d $,
    \end{itemize}
\end{frame}

%\section{Linear discriminant analysis}
%
%\section{Quadratic discriminant analysis}
%
%\section{Discriminative vs. generative}



% BibTeX slide for references. should use either acm or ieeetr style
\begin{frame}{References}
    \bibliographystyle{acm}
    % relative path may need to be updated depending on .tex file location
    \bibliography{../master_bib}
\end{frame}

\end{document}